\chapter{Inequalities and Concentration}

In Chapter 18, we saw how probability, moment, and characteristic generating
functions transform the distribution of a sum into an algebraic product.
When you know the exact distribution and can evaluate its transform in closed
form, that machinery is unmatched. But practical problems rarely hand you a
distribution clean enough to compute with. You will often know only the mean of
a quantity, perhaps its variance, or that it is built from a sum of independent
pieces --- and almost nothing else. The full distribution is a mystery. And yet
you are asked, again and again, a very practical question:

\begin{quote}
\emph{How likely is this random quantity to land far from where I expect it
to be?}
\end{quote}

This is the question of \textbf{tail probabilities} --- the probability of
landing out in the extreme fringes (the ``tails'') of a distribution, far from
where the mass clusters. The mathematical tools that answer this without
knowing the whole distribution are called \textbf{inequalities}. An
inequality trades exactness for robustness. Instead of saying ``the chance of
a big deviation is exactly $0.0137$,'' it says ``the chance is \emph{at most}
$0.02$, no matter what the distribution is, as long as it has this mean and
this variance.'' That ``no matter what'' is where the leverage comes from. It
is a guarantee that holds across an entire universe of distributions you never
had to identify.

The real payoff of these inequalities is a phenomenon called
\textbf{concentration}: when you average many independent random things, the
average stops being random in any practical sense. It \emph{concentrates}
near a single number, tighter and tighter as you add more terms, and the
probability of any meaningful deviation collapses --- often
\emph{exponentially} fast. Concentration is the reason polls work, the reason
casinos are never surprised, the reason a machine-learning model trained on a
big sample generalizes to data it has never seen. It is, with very little
exaggeration, the mathematical engine underneath all of statistics and much of
modern data science.

The chapter builds in stages. We start with Markov's inequality, the humblest
and most universal bound, relying solely on non-negativity.
Chebyshev's inequality comes next; it adds the variance and gives our first
real concentration result. From there we make the jump to \emph{exponential}
bounds --- Chernoff and Hoeffding --- which are what make concentration
genuinely powerful. Along the way we meet Jensen's inequality, the one bound in
this chapter that is about \emph{shape} rather than tails, and which quietly
underwrites half the proofs in the subject.

\chapterline

\section{The Problem: Bounding the Unknown}
\label{ch19:problem}

Let us start with a story that has no clean distribution in it at all.

A factory ships boxes of bolts. You are told one fact and one fact only: the
\emph{average} number of defective bolts per box is $2$. A buyer asks: ``What
fraction of boxes contain $10$ or more defective bolts?'' You do not know the
distribution of defects. It could be Poisson, it could be lumpy, it could be
anything. Can you say \emph{anything} useful?

It feels like you cannot. With only the mean, the distribution is wide open.
And yet there is a hard ceiling hiding in plain sight. Defect counts cannot be
negative. That single constraint --- non-negativity --- combined with the
known average, is enough to corner the answer.

\begin{intuitionbox}
The main idea is simple: \emph{If the average is small but a
value is large, then that large value cannot happen too often --- because if
it did, it would drag the average up past where we know it sits.}

Suppose, for contradiction, that more than $1/5$ of boxes had $10$ or more
defects. Those boxes alone would contribute more than $\tfrac15 \times 10 = 2$
to the average. Every other box contributes a non-negative amount on top of
that. So the overall average would have to exceed $2$. But the average
\emph{is} $2$. Contradiction. Therefore at most $1/5$ of boxes can have $10$
or more defects.
\end{intuitionbox}

That little argument is the entire content of Markov's inequality. No
distribution required. Just ``values are non-negative'' and ``here is the
mean.'' Let us make it a theorem.

\section{Markov's Inequality: Everything From Non-Negativity}
\label{ch19:markov}

\begin{theorem}[Markov's inequality]
\label{ch19:markov-thm}
Let $X$ be a random variable with $X \ge 0$, and let $a > 0$. Then
\[
  \P(X \ge a) \;\le\; \frac{\E[X]}{a}.
\]
\end{theorem}

\begin{whybox}
The proof is the sand-on-a-seesaw picture turned into one line. The trick is to
build a crude ``staircase'' function that lies underneath $X$ but is easy to
average.

Fix $a>0$ and look at the indicator $\ind\set{X \ge a}$, which is $1$ when
$X \ge a$ and $0$ otherwise. We claim the simple inequality
\[
  a \cdot \ind\set{X \ge a} \;\le\; X
  \qquad \text{(always, for every outcome).}
\]
Check the two cases. If $X \ge a$, the left side is $a$ and the right side is
$X \ge a$, so the inequality holds. If $X < a$, the left side is $0$ and the
right side is $X \ge 0$ (here we use non-negativity!), so it holds again. The
inequality is true outcome by outcome, so it survives taking expectations:
\[
  a\,\E\!\left[\ind\set{X\ge a}\right] \;\le\; \E[X].
\]
But the expectation of an indicator is just the probability of the event,
$\E[\ind\set{X\ge a}] = \P(X\ge a)$. Divide by $a>0$ and you are done.
Figure~\ref{ch19:fig-markov} shows the staircase-under-the-line picture that
the algebra is tracking.
\end{whybox}

\begin{figure}[H]
\centering
\begin{tikzpicture}[scale=0.95]
  % Panel 1: Markov linear bound
  \begin{scope}[shift={(0,0)}]
    \draw[draw=gray!30, fill=gray!3, rounded corners=4pt] (-0.8,-1.2) rectangle (5.8,4.2);
    \node[anchor=north, font=\small\bfseries, densityred!90!black] at (2.5,4.05) {1. Markov: $a\,\ind\{X \ge a\} \le X$};

    \draw[->, thick, gray!70] (-0.2,0) -- (4.8,0) node[right, font=\scriptsize] {$x$};
    \draw[->, thick, gray!70] (0,-0.2) -- (0,3.2) node[above, font=\scriptsize] {$y$};

    \draw[densityred, very thick] (0,0) -- (3.0,3.0) node[right, font=\scriptsize\bfseries] {$y = x$};

    \coordinate (A) at (1.8, 1.8);
    \draw[sourceblue!90!black, very thick] (0,0) -- (1.8,0);
    \draw[sourceblue!90!black, very thick] (1.8,1.8) -- (4.2,1.8)
      node[midway, above=2pt, font=\scriptsize\bfseries] {$y = a\,\ind\{x \ge a\}$};
    \draw[dashed, gray] (1.8,0) -- (1.8,1.8);

    \draw[gray!70] (1.8,0.08) -- (1.8,-0.08) node[below, font=\scriptsize] {$a$};
    \draw[gray!70] (0.08,1.8) -- (-0.08,1.8) node[left, font=\scriptsize] {$a$};

    \fill[sourceblue!18] (1.8,0) rectangle (4.2,1.8);

    \node[align=center, font=\scriptsize, draw=sourceblue!70, fill=white, rounded corners=2pt, inner sep=2pt]
      at (2.3, -0.6) {Step lies below line $\implies a\,\P(X \ge a) \le \E[X]$};
  \end{scope}

  % Transmission arrow
  \draw[->, >=stealth, very thick, structurepurple] (6.0, 1.3) -- (7.1, 1.3)
    node[midway, above, font=\scriptsize\bfseries, align=center] {Apply to\\$(X-\mu)^2$};

  % Panel 2: Chebyshev quadratic bowl bound
  \begin{scope}[shift={(10.1,0)}]
    \draw[draw=gray!30, fill=gray!3, rounded corners=4pt] (-2.8,-1.2) rectangle (4.0,4.2);
    \node[anchor=north, font=\small\bfseries, quotientorange!90!black] at (0.6,4.05) {2. Chebyshev: $k^2\,\ind\{|u|\ge k\} \le u^2$};

    \draw[->, thick, gray!70] (-2.5,0) -- (3.0,0) node[right, font=\scriptsize] {$u = x - \mu$};
    \draw[->, thick, gray!70] (0,-0.2) -- (0,3.2) node[above, font=\scriptsize] {$y$};

    \draw[densityred, very thick, domain=-1.65:1.65, samples=60]
      plot (\x, {\x*\x}) node[above right=-2pt, font=\scriptsize\bfseries] {$y = u^2$};

    \coordinate (K1) at (-1.1, 1.21);
    \coordinate (K2) at (1.1, 1.21);

    \draw[sourceblue!90!black, very thick] (-2.4,1.21) -- (-1.1,1.21);
    \draw[sourceblue!90!black, very thick] (-1.1,0) -- (1.1,0);
    \draw[sourceblue!90!black, very thick] (1.1,1.21) -- (2.4,1.21)
      node[right, font=\scriptsize\bfseries] {$y = k^2\,\ind\{|u|\ge k\}$};

    \draw[dashed, gray] (-1.1,0) -- (-1.1,1.21);
    \draw[dashed, gray] (1.1,0) -- (1.1,1.21);

    \draw[gray!70] (-1.1,0.08) -- (-1.1,-0.08) node[below, font=\scriptsize] {$-k$};
    \draw[gray!70] (1.1,0.08) -- (1.1,-0.08) node[below, font=\scriptsize] {$k$};
    \draw[gray!70] (0.08,1.21) -- (-0.08,1.21) node[left, font=\scriptsize] {$k^2$};

    \fill[sourceblue!18] (-2.4,0) rectangle (-1.1,1.21);
    \fill[sourceblue!18] (1.1,0) rectangle (2.4,1.21);

    \node[align=center, font=\scriptsize, draw=quotientorange!80!black, fill=white, rounded corners=2pt, inner sep=2pt]
      at (0.3, -0.6) {Bowl dominates step $\implies k^2\,\P(|X-\mu|\ge k) \le \sigma^2$};
  \end{scope}
\end{tikzpicture}
\caption{The geometric mechanism of Markov and Chebyshev bounds. (Left) The step function $a\,\ind\{X \ge a\}$ is pointwise bounded above by the line $y=x$ for $x \ge 0$, forcing the expected step area $a\,\P(X \ge a)$ to lie below the mean $\E[X]$. (Right) Applying the exact same staircase logic to the non-negative quadratic bowl $u^2 = (X-\mu)^2$ yields a two-sided step of height $k^2$ bounded by $\E[(X-\mu)^2] = \sigma^2$, delivering Chebyshev's inequality $\P(|X-\mu| \ge k) \le \sigma^2/k^2$.}
\label{ch19:fig-markov}
\end{figure}

\begin{happeningbox}
Markov's inequality looks almost too weak to be useful --- and used naively,
it often is. Its real power is that it is a \emph{template}. The proof never
cared that $X$ was the original quantity of interest. If you can find
\emph{any} clever non-negative function $g(X) \ge 0$ that blows up when $X$
strays, you can apply Markov to $g(X)$ and get a bound on $X$ straying. Feed
it $g(X) = (X-\mu)^2$ and out pops Chebyshev. Feed it
$g(X) = e^{tX}$ and out pops the Chernoff bound, the most powerful tail
inequality in this book. Every concentration result we prove is an application
of Markov's inequality to an appropriately chosen function.
\end{happeningbox}

\begin{example}[The bolt boxes, finished]
\label{ch19:ex-bolts}
With $X = $ defects per box, $\E[X]=2$, and $a=10$:
\[
  \P(X \ge 10) \;\le\; \frac{\E[X]}{10} = \frac{2}{10} = \frac15.
\]
At most $20\%$ of boxes have $10$ or more defective bolts --- guaranteed,
with no knowledge of the distribution. If the truth happened to be Poisson
with mean $2$, the real probability is about $0.0000465$ (that is,
$4.6\times 10^{-5}$), so Markov is wildly loose here. That looseness is the price of universality: the bound has to hold
even for the worst-case distribution, which in this case is a nasty one that
puts a tiny sliver of mass exactly at $10$ and the rest at $0$.
\end{example}

\begin{walkbox}
A city's average household income is \$70,000. Can 80\% of households make
\$100,000 or more? Check in your head in five seconds:
$\P(X \ge 100{,}000) \le 70{,}000 / 100{,}000 = 70\%$. Because $80\% > 70\%$,
it is mathematically impossible (assuming non-negative incomes). Those $80\%$
alone would average at least $0.80 \times \$100{,}000 = \$80{,}000$, exceeding
the city average before the remaining $20\%$ of households are even counted.
\end{walkbox}

\begin{trapbox}
Markov needs $X \ge 0$. Forgetting this is the classic blunder. If $X$ can be
negative, the staircase argument breaks at the step ``$X < a$ implies
$X \ge 0$,'' and the bound is simply false. To bound a two-sided quantity like
$X-\mu$, do not apply Markov to it directly --- apply Markov to a non-negative
\emph{function} of it, such as $(X-\mu)^2$. That detour is exactly how we get
Chebyshev next.
\end{trapbox}

\begin{marginnotebox}
The inequality carries the name of Andrey Markov (1856--1922), the St.\
Petersburg mathematician who also introduced Markov chains. Yet the underlying
insight --- that a non-negative quantity with a small mean cannot frequently
take giant values --- was used decades earlier by his teacher Pafnuty Chebyshev
to prove the variance-based bound we meet next. Markov is the raw engine;
Chebyshev is the polished theorem built upon it.
\end{marginnotebox}

\section{Chebyshev's Inequality: Bringing In the Variance}
\label{ch19:chebyshev}

Markov used only the mean. But if you also know the \emph{variance} --- how
spread out $X$ is --- you should be able to say much more about how far $X$
wanders. Variance is literally a measure of typical squared distance from the
mean, so it is exactly the right currency for buying a tail bound.

The move is the one we previewed above: apply Markov not to $X$, but to
the non-negative quantity $(X-\mu)^2$.

\begin{theorem}[Chebyshev's inequality]
\label{ch19:cheb-thm}
Let $X$ have mean $\mu = \E[X]$ and finite variance $\sigma^2 = \Var(X)$. Then
for any $k > 0$,
\[
  \P\bigl(\abs{X-\mu} \ge k\bigr) \;\le\; \frac{\sigma^2}{k^2}.
\]
Equivalently, writing the deviation in units of standard deviations
($k = c\,\sigma$),
\[
  \P\bigl(\abs{X-\mu} \ge c\,\sigma\bigr) \;\le\; \frac{1}{c^2}.
\]
\end{theorem}

\begin{whybox}
Let $Y = (X-\mu)^2$. This $Y$ is non-negative, so Markov applies. Its mean is
the variance: $\E[Y] = \E[(X-\mu)^2] = \sigma^2$. Now notice that the events
$\set{\abs{X-\mu}\ge k}$ and $\set{Y \ge k^2}$ are the \emph{same event} ---
$\abs{X-\mu}\ge k$ exactly when $(X-\mu)^2 \ge k^2$. So
\[
  \P\bigl(\abs{X-\mu}\ge k\bigr)
  = \P\bigl(Y \ge k^2\bigr)
  \;\overset{\text{Markov}}{\le}\;
  \frac{\E[Y]}{k^2}
  = \frac{\sigma^2}{k^2}.
\]
That is the whole proof. Squaring turned a two-sided deviation into a
one-sided statement about a non-negative variable, and Markov finished the
job.
\end{whybox}

The ``standard deviations'' form is the one worth memorizing, because it is
distribution-free and strikingly clean:

\begin{itemize}
  \item At least $1 - 1/2^2 = 75\%$ of any distribution's mass lies within
        $2$ standard deviations of the mean.
  \item At least $1 - 1/3^2 \approx 88.9\%$ lies within $3$ standard
        deviations.
  \item At least $1 - 1/4^2 \approx 93.75\%$ lies within $4$.
\end{itemize}

No bell-curve assumption, no symmetry, nothing. Any random variable with a
finite variance obeys these; Figure~\ref{ch19:fig-cheb} plots the ceiling
$1/c^2$ against the much smaller Gaussian tail. (The Normal distribution does far better ---
about $95\%$ within two standard deviations and $99.7\%$ within three --- but
that requires \emph{knowing} you have a Normal. Chebyshev is the floor that
holds for everyone.)

\begin{figure}[H]
\centering
\begin{tikzpicture}
\begin{axis}[
  width=13cm, height=6cm,
  xmin=1, xmax=5, ymin=0, ymax=1.05,
  xlabel={$c$ = deviation in standard deviations},
  ylabel={tail probability bound},
  domain=1:5, samples=120,
  legend pos=north east,
  grid=both, grid style={gray!18},
]
\addplot[densityred, thick] {1/(x^2)};
\addlegendentry{Chebyshev bound $1/c^2$}
\addplot[sourceblue, thick, dashed] {exp(-x^2/2)};
\addlegendentry{Gaussian tail bound $e^{-c^2/2}$}
\end{axis}
\end{tikzpicture}
\caption{The Chebyshev ceiling $1/c^2$ (red) bounds the probability of being
at least $c$ standard deviations from the mean, for \emph{any} distribution.
For comparison the Gaussian tail bound $e^{-c^2/2}$ (blue dashed), which
controls the Normal, is far smaller --- Chebyshev pays for its universality
with looseness, but it holds without any distributional assumptions.}
\label{ch19:fig-cheb}
\end{figure}

\begin{example}[Polling, the crude version]
\label{ch19:ex-poll}
You poll $n$ people and let $\hat p$ be the fraction who say ``yes.'' Each
response is $\Bern(p)$, so $\hat p$ has mean $p$ and variance
$p(1-p)/n \le 1/(4n)$ (the variance of a Bernoulli is largest at $p=1/2$).
Chebyshev says
\[
  \P\bigl(\abs{\hat p - p} \ge \eps\bigr)
  \;\le\; \frac{\Var(\hat p)}{\eps^2}
  \;\le\; \frac{1}{4n\eps^2}.
\]
To be within $\eps = 0.03$ of the truth with probability at least $95\%$, set
$1/(4n\eps^2) \le 0.05$, giving $n \ge 1/(4\cdot 0.05 \cdot 0.0009) \approx
5556$. So a few thousand respondents pin down a national opinion to within
three points --- and this is the \emph{pessimistic} Chebyshev count. The
exponential bounds later in the chapter will sharpen it, but the headline is
already here: accuracy depends on $n$, not on the size of the population.
\end{example}

\begin{happeningbox}
Look at what Chebyshev does to an \emph{average}. Let
$\bar X_n = \frac1n\sum_{i=1}^n X_i$ be the average of $n$ independent copies,
each with mean $\mu$ and variance $\sigma^2$. Variances of independent things
add, and the $1/n$ out front gets squared, so
$\Var(\bar X_n) = \sigma^2/n$. Plug into Chebyshev:
\[
  \P\bigl(\abs{\bar X_n - \mu} \ge \eps\bigr)
  \;\le\; \frac{\sigma^2}{n\,\eps^2}
  \;\xrightarrow[n\to\infty]{}\; 0.
\]
The right-hand side melts to zero as $n$ grows. The average concentrates on
$\mu$. This single line is a complete proof of the \textbf{Weak Law of Large
Numbers} --- the formal statement that sample averages converge in
probability to the true mean. We will revisit it in its own chapter, but you
have just watched it fall out of one inequality. \emph{This} is concentration:
the spread of the average shrinks to zero at an explicit, quantifiable rate.
\end{happeningbox}

\section{What ``Concentration'' Really Means}
\label{ch19:concentration}

Pause to appreciate the shift in worldview. A single fair coin flip is maximally
uncertain --- you cannot predict it at all. But the \emph{average} of a million
fair coin flips is almost perfectly predictable: it will be $0.5$ to within a
whisker, essentially every time. The randomness has not disappeared; it has
been \emph{diluted}. Each flip still surprises you, but their collective
average has nowhere to hide. Figure~\ref{ch19:fig-concentration} shows the
distribution of the sample average tightening around $0.5$ as $n$ grows.

\begin{figure}[H]
\centering
\begin{tikzpicture}
\begin{axis}[
  width=13cm, height=6cm,
  xmin=0, xmax=1,
  ymin=0, ymax=12,
  xlabel={value of the sample average $\bar X_n$},
  ylabel={density (schematic)},
  domain=0:1, samples=160,
  legend pos=north west,
  ytick=\empty,
]
\addplot[sourceblue, thick] {1/(0.18*sqrt(2*pi))*exp(-((x-0.5)^2)/(2*0.18^2))};
\addlegendentry{$n=10$}
\addplot[quotientorange, thick] {1/(0.08*sqrt(2*pi))*exp(-((x-0.5)^2)/(2*0.08^2))};
\addlegendentry{$n=50$}
\addplot[densityred, thick] {1/(0.035*sqrt(2*pi))*exp(-((x-0.5)^2)/(2*0.035^2))};
\addlegendentry{$n=250$}
\draw[gray, dashed] (axis cs:0.5,0) -- (axis cs:0.5,12);
\end{axis}
\end{tikzpicture}
\caption{Concentration in action. As the number of averaged coin flips $n$
grows, the distribution of the sample average $\bar X_n$ piles up ever more
tightly around the true mean $0.5$. The spread shrinks like $1/\sqrt n$; the
\emph{tail probability} of any fixed deviation shrinks far faster, as the next
section shows.}
\label{ch19:fig-concentration}
\end{figure}

\begin{aestheticbox}
Chebyshev's $\sigma^2/(n\eps^2)$ tells us the tail probability shrinks at least
as fast as $1/n$. That is genuinely useful, and for a long time it was the best
general tool available. But $1/n$ is, for many purposes, embarrassingly slow.
To drive a failure probability down to one in a million, $1/n$ decay demands
about a million samples. The remarkable discovery of the twentieth century is
that for \emph{sums of independent bounded variables}, the true decay is
\emph{exponential} in $n$: the tail behaves like $e^{-cn}$, not $1/n$. One in a
million then costs only a few dozen samples' worth of exponent. The leap from
polynomial to exponential decay is the difference between statistics that
``eventually works'' and statistics that works \emph{decisively}. The tool
that captures it is the Chernoff method.
\end{aestheticbox}

\section{The Chernoff Method: Exponential Tails}
\label{ch19:chernoff}

Why is Chebyshev slow? Because $(X-\mu)^2$ grows only quadratically, so a
single far-out outlier is not penalized very harshly inside the expectation.
To get an exponential tail bound we need an exponentially harsh penalty. So we
feed Markov's inequality the most aggressive non-negative function we have:
the exponential $e^{tX}$.

Here is the trick in full. For any $t > 0$, the event $\set{X \ge a}$ is
identical to $\set{e^{tX} \ge e^{ta}}$, because $e^{t\,\cdot}$ is increasing.
Apply Markov to the non-negative variable $e^{tX}$:
\[
  \P(X \ge a) \;=\; \P\bigl(e^{tX} \ge e^{ta}\bigr)
  \;\le\; \frac{\E[e^{tX}]}{e^{ta}}
  \;=\; e^{-ta}\,M_X(t),
\]
where $M_X(t) = \E[e^{tX}]$ is the \textbf{moment generating function}
introduced in Chapter 18 --- the expectation that tracks all moments of $X$
and converts sums of independent terms into products. This holds for
\emph{every} $t > 0$ simultaneously, so we are free to choose the $t$ that
makes the bound tightest:
\[
  \boxed{\;\P(X \ge a) \;\le\; \min_{t>0}\; e^{-ta}\,M_X(t).\;}
\]
This is the \textbf{Chernoff bound}. The minimization over $t$ is what gives it
its power.

\begin{whybox}
Why does optimizing over $t$ produce exponential decay? Think of $e^{-ta}M_X(t)$
as a family of bounds, one per $t$. A small $t$ barely penalizes large values,
giving a weak bound; a huge $t$ penalizes them ferociously but blows up
$M_X(t)$ itself. Somewhere in between is a sweet spot where the two effects
balance, and at that sweet spot the exponent $-ta$ wins out, scaling linearly
in $a$ inside an exponential. Because $M_X(t)$ for a \emph{sum} of $n$
independent variables factorizes --- $M_{S_n}(t) = M_X(t)^n$ --- the whole
bound becomes $\bigl(e^{-ta/n}M_X(t)\bigr)^n$, an $n$-th power. A number less
than one raised to the $n$ is exponential decay in $n$. That factorization,
inherited from independence, is the engine. Chebyshev never exploited it
because squaring does not turn products into powers the way exponentials do.
Figure~\ref{ch19:fig-chernoff-convex} illustrates this convex sweet spot
geometrically.
\end{whybox}

\begin{figure}[H]
\centering
\begin{tikzpicture}
\begin{axis}[
  width=12cm, height=5.8cm,
  xmin=0, xmax=2.5, ymin=0.85, ymax=1.2,
  xlabel={$t$ (tilt parameter)}, ylabel={$g(t) = e^{-ta} M(t)$},
  domain=0:2.2, samples=100,
  legend pos=north east,
  title={The Chernoff optimization bowl for a fair coin: $a = 0.7$, $\mu = 0.5$},
  title style={font=\small},
]
\addplot[sourceblue, very thick] {0.5*(exp(-0.7*x) + exp(0.3*x))};
\addlegendentry{$g(t) = e^{-0.7 t} M(t)$}
\draw[densityred, dashed] (axis cs:0.847,0.85) -- (axis cs:0.847,0.921);
\fill[densityred] (axis cs:0.847,0.9206) circle (2.2pt);
\node (spot) [font=\scriptsize\bfseries, densityred!90!black, fill=white, inner sep=2pt] at (axis cs:0.847,1.06) {sweet spot $t^\star \approx 0.85$, $g(t^\star) \approx 0.921$};
\draw[densityred, ->, >=stealth, thick, shorten <=6pt] (spot.south) -- (axis cs:0.847,0.932);
\node[above right=2pt, font=\scriptsize, gray!60!black] at (axis cs:0,1.0) {$g(0) = 1$};
\end{axis}
\end{tikzpicture}
\caption{The Chernoff optimization mechanism. For any target deviation $a > \E[X]$, the bound $g(t) = e^{-ta}M(t)$ begins at $g(0) = 1$ with an initial downward slope of $\E[X] - a < 0$. As $t$ increases, the exponential penalty begins to dominate, curving the function upward into a convex valley before the moment generating function blows up. The lowest point at $t^\star$ gives the sharpest bound $g(t^\star) < 1$. When averaging $n$ independent trials, this single-trial factor is raised to the $n$-th power: $(g(t^\star))^n$, collapsing exponentially to zero.}
\label{ch19:fig-chernoff-convex}
\end{figure}

\begin{example}[Chernoff for a sum of fair coins]
\label{ch19:ex-chernoff-coins}
Let $S_n = X_1 + \cdots + X_n$ be the number of heads in $n$ fair flips, so
$\E[S_n] = n/2$. Each $X_i$ has MGF
$\E[e^{tX_i}] = \tfrac12(1 + e^{t})$, and by independence
$M_{S_n}(t) = \bigl(\tfrac12(1+e^t)\bigr)^n$. The Chernoff bound on the chance
of seeing at least $(\tfrac12 + \delta)n$ heads is
\[
  \P\!\left(S_n \ge \bigl(\tfrac12+\delta\bigr)n\right)
  \;\le\;
  \min_{t>0}\, e^{-t(\frac12+\delta)n}\Bigl(\tfrac12(1+e^t)\Bigr)^n .
\]
To minimize the exponent, we can use the clean upper bound
$\tfrac12(1+e^t) \le e^{t/2 + t^2/8}$ (a consequence of the convexity of the
exponential, formalised as Hoeffding's lemma below). With this bound, the
per-trial factor becomes
\[
  e^{-t(\frac12+\delta)}\,\E[e^{tX_i}] \;\le\; e^{-t\delta + t^2/8}.
\]
Differentiating the exponent $-t\delta + t^2/8$ with respect to $t$ gives
$-\delta + t/4 = 0$, so the optimal choice is $t^\star = 4\delta$. Substituting
$t^\star = 4\delta$ back in yields $-4\delta^2 + (4\delta)^2/8 = -2\delta^2$,
delivering the clean exponential
\[
  \P\!\left(S_n \ge \bigl(\tfrac12+\delta\bigr)n\right)
  \;\le\; e^{-2\delta^2 n}.
\]
Compare to Chebyshev, which here gives only
$\P(\abs{S_n - n/2}\ge \delta n) \le 1/(4\delta^2 n)$, a measly $1/n$. For a
$10\%$ excess ($\delta = 0.1$) with $n = 1000$ flips, Chebyshev allows up to
$1/40 = 2.5\%$; Chernoff caps it at $e^{-20} \approx 2\times 10^{-9}$. This
illustrates the dramatic difference between polynomial and exponential decay,
shown on a logarithmic scale in Figure~\ref{ch19:fig-chernoff-vs-cheb}.
\end{example}

\begin{figure}[H]
\centering
\begin{tikzpicture}
\begin{axis}[
  width=13cm, height=6.2cm,
  xmin=0, xmax=1000,
  ymin=1e-12, ymax=1,
  ymode=log,
  xlabel={number of samples $n$},
  ylabel={upper bound on tail probability},
  legend style={at={(0.05, 0.08)}, anchor=south west},
  domain=10:1000, samples=80,
  grid=both, grid style={gray!18},
]
\addplot[densityred, thick] {1/(4*0.1^2*x)};
\addlegendentry{Chebyshev $\sim 1/n$}
\addplot[sourceblue, thick] {exp(-2*0.1^2*x)};
\addlegendentry{Chernoff $\sim e^{-2\delta^2 n}$}
\end{axis}
\end{tikzpicture}
\caption{Tail bound on $\P(S_n \ge (\tfrac12+\delta)n)$ for $\delta=0.1$, on a
log scale. Chebyshev (red) crawls down like $1/n$; Chernoff (blue) plunges
exponentially. By $n=1000$ they differ by many orders of magnitude. This gap
is why concentration is taken seriously.}
\label{ch19:fig-chernoff-vs-cheb}
\end{figure}

\begin{historybox}
The story of concentration began with a struggle over rigor in 19th-century
St.\ Petersburg. In the 1860s, \textbf{Pafnuty Chebyshev} wanted to prove that
averages settle down to their expected values (the Law of Large Numbers) for
arbitrary distributions. The existing proofs, dating back to Jacob Bernoulli
and Sim\'eon Denis Poisson, relied on Stirling's approximation formula for
factorials --- an approach locked to coin-tossing binomials and helpless for
general variables. Chebyshev had only basic algebra and the definitions of mean
and variance. His breakthrough was seeing that squaring $(X-\mu)^2$ forces all
deviations to be positive, allowing variance to serve as a universal ceiling.
His student \textbf{Andrey Markov} later isolated the raw non-negative
inequality that underpins Chebyshev's step.

Yet for decades, Western applied statisticians greeted Chebyshev's inequality
with skepticism, dismissing it as ``mathematically elegant but practically
useless.'' Because its tail bound decays only like $1/k^2$, it guarantees that
at least $75\%$ of any distribution lies within two standard deviations,
whereas for bell-shaped empirical data the truth is closer to $95\%$.
To practitioners, $1/k^2$ felt like an absurdly loose price to pay for
distribution-free safety.

The breakthrough came from changing the currency. In the 1920s, the Russian
mathematician \textbf{Sergei Bernstein} realized that if you replace the
quadratic bowl with the exponential function $e^{tX}$, independence converts
the expectation of a sum into a product of expectations, trading polynomial
crawl for exponential collapse. In 1952, \textbf{Herman Chernoff} published the
modern optimization recipe, generously crediting his Stanford colleague
\textbf{Herman Rubin} with the insight to minimize over the tilt parameter $t$.
A decade later, in 1963, \textbf{Wassily Hoeffding} proved that if variables
are simply trapped inside a bounded interval $[a,b]$, their exponential moment
is controlled by a clean quadratic exponent, freeing practitioners from needing
to calculate the moment generating function at all.

Concentration has since moved far beyond sums of independent numbers. In the
late 20th century, \textbf{Vitali Milman} and \textbf{Michel Talagrand}
generalized concentration to geometric functions on high-dimensional spaces,
discovering that in high dimensions, almost all the volume of a sphere or cube
is concentrated in an infinitesimally thin shell near its equator. This
``concentration of measure'' became the mathematical foundation of randomized
algorithms in computer science and \textbf{PAC learning} (``probably
approximately correct'' learning) in machine learning, proving that statistical
models can generalize reliably from finite training samples.

Active research frontiers remain open. First, classical concentration relies
heavily on independence; developing equally sharp, non-asymptotic bounds for
\emph{dependent} processes --- such as time series, Markov processes, or
physical systems with local interactions --- remains an active research area.
Second, in high-dimensional statistics and network analysis, researchers study
random \emph{matrices} rather than random scalars; proving concentration for
matrix eigenvalues requires overcoming the non-commutativity of matrix products
($AB \ne BA$). Finally, for \emph{heavy-tailed} data in finance and physics
where the moment generating function $M_X(t)$ diverges, researchers develop
truncation and polynomial concentration techniques where exponential methods
fail.
\end{historybox}

\section{Hoeffding's Inequality: Concentration You Can Use}
\label{ch19:hoeffding}

The Chernoff recipe is powerful but demands the MGF, which you may not know.
\textbf{Hoeffding's inequality} packages the recipe into a ready-to-use form
that needs only one mild assumption: each variable is \emph{bounded}. This is
the concentration inequality you will reach for most often in practice.

\begin{theorem}[Hoeffding's inequality]
\label{ch19:hoeff-thm}
Let $X_1, \dots, X_n$ be independent with $a_i \le X_i \le b_i$ almost surely.
Let $S_n = \sum_{i=1}^n X_i$. Then for any $\eps > 0$,
\[
  \P\bigl(S_n - \E[S_n] \ge \eps\bigr)
  \;\le\;
  \exp\!\left(-\frac{2\eps^2}{\sum_{i=1}^n (b_i - a_i)^2}\right),
\]
and the same bound holds for the lower tail $\P(\E[S_n] - S_n \ge \eps)$. In
the common case where each $X_i \in [a,b]$ and we look at the
\emph{average} $\bar X_n = S_n/n$,
\[
  \P\bigl(\abs{\bar X_n - \mu} \ge \eps\bigr)
  \;\le\;
  2\exp\!\left(-\frac{2 n \eps^2}{(b-a)^2}\right).
\]
\end{theorem}

\begin{whybox}
Hoeffding is Chernoff plus one technical lemma. The Chernoff step gives, for a
centered bounded variable $Y_i = X_i - \E[X_i]$,
\[
  \P(S_n - \E[S_n] \ge \eps)
  \le e^{-t\eps}\prod_{i=1}^n \E\!\bigl[e^{tY_i}\bigr],
\]
using independence to split the expectation of the product into a product of
expectations. The missing piece is a bound on each factor $\E[e^{tY_i}]$.
\textbf{Hoeffding's lemma} supplies it: if $Y$ has mean $0$ and lies in
$[a,b]$, then $\E[e^{tY}] \le e^{t^2(b-a)^2/8}$. (Intuitively, among all
mean-zero variables confined to an interval, the one with the fattest
exponential moment is the two-point distribution sitting at the endpoints,
and a short convexity argument pins down the $(b-a)^2/8$ constant.) Plugging
in,
\[
  \P(S_n - \E[S_n] \ge \eps)
  \le \exp\!\left(-t\eps + \tfrac{t^2}{8}\sum_i (b_i-a_i)^2\right).
\]
This is a parabola in $t$ inside the exponent; minimizing over $t$ (vertex at
$t^\star = 4\eps/\sum_i (b_i-a_i)^2$) gives exactly the stated
$-2\eps^2/\sum_i(b_i-a_i)^2$. The whole proof is: Chernoff splits, Hoeffding's
lemma bounds each piece, calculus finds the best $t$.
\end{whybox}

\begin{happeningbox}
Examine the exponent in the bound for the sample mean, $-2n\eps^2/(b-a)^2$. Two features
matter enormously.

First, $n$ sits in the exponent. Double your sample size and you do not halve
the failure probability --- you \emph{square} it. This is the exponential
concentration that Chebyshev's $1/n$ could only dream of.

Second, the bound depends on the range $(b-a)$ but \emph{not} on the variance,
the shape, or anything else about the distribution. Bounded plus independent is
the entire price of admission. That distribution-freeness, married to
exponential decay, is precisely the combination that statistical learning
theory needs --- which is why Hoeffding is the first inequality in every
machine-learning theory course.
\end{happeningbox}

\begin{example}[Polling, the sharp version]
\label{ch19:ex-poll-sharp}
Return to the poll. Each response $X_i \in [0,1]$, so $b-a = 1$, and
$\bar X_n = \hat p$. Hoeffding gives
\[
  \P\bigl(\abs{\hat p - p} \ge \eps\bigr)
  \;\le\; 2 e^{-2n\eps^2}.
\]
For $\eps = 0.03$ and a $95\%$ guarantee we need $2e^{-2n(0.03)^2} \le 0.05$,
i.e.\ $e^{-0.0018\,n} \le 0.025$, giving
$n \ge \ln(40)/0.0018 \approx 2050$. Recall Chebyshev demanded about
$5556$ for the same guarantee. Same data, same assumptions, sharper tool ---
less than half the sample size. Figure~\ref{ch19:fig-hoeffding} traces how the
guaranteed half-width shrinks with $n$. And unlike Chebyshev, Hoeffding hands you a
confidence \emph{interval} of half-width $\eps = \sqrt{\ln(2/\alpha)/(2n)}$ at
confidence $1-\alpha$, the Hoeffding bound used throughout applied statistics.
\end{example}

\begin{figure}[H]
\centering
\begin{tikzpicture}
\begin{axis}[
  width=13cm, height=6cm,
  xmin=0, xmax=4000,
  ymin=0, ymax=0.12,
  ytick={0,0.05,0.10},
  scaled y ticks=false,
  yticklabel style={/pgf/number format/fixed, /pgf/number format/precision=2},
  xlabel={number of samples $n$},
  ylabel={Hoeffding half-width $\eps$},
  domain=50:4000, samples=120,
  legend pos=north east,
  grid=both, grid style={gray!18},
]
\addplot[sourceblue, thick] {sqrt(ln(2/0.05)/(2*x))};
\addlegendentry{$95\%$ confidence ($\alpha=0.05$)}
\addplot[densityred, thick, dashed] {sqrt(ln(2/0.01)/(2*x))};
\addlegendentry{$99\%$ confidence ($\alpha=0.01$)}
\draw[gray, dashed] (axis cs:0,0.03) -- (axis cs:4000,0.03);
\end{axis}
\end{tikzpicture}
\caption{The Hoeffding confidence half-width $\eps=\sqrt{\ln(2/\alpha)/(2n)}$
for a bounded average on $[0,1]$, as a function of sample size $n$. It shrinks
like $1/\sqrt n$: quadrupling the sample halves the width. The horizontal line
marks $\eps=0.03$; the $95\%$ curve crosses it near $n\approx 2050$, matching
the poll calculation. Demanding higher confidence ($99\%$, red dashed) widens
the band only mildly, because confidence enters through the slow logarithm
$\ln(2/\alpha)$.}
\label{ch19:fig-hoeffding}
\end{figure}

\begin{applicationbox}
\textbf{Why concentration is the backbone of statistics and machine learning.}

Every estimate you ever compute --- a sample mean, a regression coefficient,
a model's accuracy on a test set --- is a function of random data. The entire
enterprise of statistics rests on a single hope: that this function does not
swing wildly with the luck of the sample, but \emph{concentrates} near the
true value. Concentration inequalities turn that hope into a theorem with an
explicit, computable failure probability.

In supervised machine learning the stakes are sharper still. You train a model
to minimize error on a finite \emph{training} sample, but you care about its
error on the unseen population --- its \emph{generalization} error. The two
differ by exactly the kind of deviation Hoeffding controls. For a single fixed
hypothesis $h$, the training error is an average of $n$ bounded $\set{0,1}$
mistakes, so Hoeffding gives
\[
  \P\bigl(\abs{\widehat{\mathrm{err}}(h) - \mathrm{err}(h)} \ge \eps\bigr)
  \le 2 e^{-2n\eps^2}.
\]
To cover a whole hypothesis class of size $M$ (so we may safely \emph{choose}
the best $h$ after seeing the data) one applies the union bound, paying a
factor $M$ that becomes a mild additive $\ln M$ inside the exponent's budget:
the sample size needed grows only like $\log M$. This is the seed of
\textbf{PAC learning} (``probably approximately correct'') and of VC theory,
which extend the argument to infinite classes via cleverer counting. Strip
away the formalism and the message is simple: \emph{a model generalizes
because averages concentrate, and we can bound the sample size that makes the
guarantee bite.} Concentration is what turns good performance on a finite
training sample into a guarantee on unseen data.
\end{applicationbox}

\begin{connectionbox}
The Chernoff and Hoeffding proofs lean entirely on the moment generating
function $M_X(t) = \E[e^{tX}]$ from the chapter on generating functions, and
the optimization step is plain single-variable calculus --- find the $t$ where
the derivative of the exponent vanishes. The deeper link is to
\emph{large-deviations theory}: the limit $-\frac1n \log \P(\bar X_n \ge a)$
converges to the \textbf{Legendre transform} of $\log M_X(t)$, a construction
borrowed straight from convex analysis and classical mechanics. The Chernoff
bound is the finite-$n$ shadow of that exact asymptotic rate. If you enjoyed
the Legendre transform in the companion calculus book, you are meeting its
probabilistic counterpart here.
\end{connectionbox}

\section{Jensen's Inequality: The Bound About Shape}
\label{ch19:jensen}

The inequalities so far were about \emph{tails} --- how often a variable
strays. Jensen's inequality is about something different and equally
fundamental: how \emph{averaging interacts with curvature}. It answers the
question, ``is the average of a function the function of the average?'' ---
and the answer, almost always, is no, with a definite direction to the error.

\begin{intuitionbox}
Picture a function that bends upward, like $g(x) = x^2$ --- a \emph{convex}
function, holding water. Take two input points and look at the chord
connecting the corresponding points on the curve. Because the curve bends away
below the chord, the chord sits \emph{above} the curve. Now read this in the
language of averages: the height of the chord at the midpoint is the
\emph{average of the function values}, while the curve's height at the midpoint
is the \emph{function of the average input}. Chord above curve means
\[
  \text{average of } g \;\ge\; g(\text{of the average}).
\]
That is Jensen's inequality, which extends the idea that ``convex curves lie
below their chords,'' generalized from two points to an entire probability
distribution of points.
\end{intuitionbox}

\begin{theorem}[Jensen's inequality]
\label{ch19:jensen-thm}
Let $g$ be a convex function and $X$ a random variable with finite mean. Then
\[
  g\bigl(\E[X]\bigr) \;\le\; \E\bigl[g(X)\bigr].
\]
If $g$ is concave, the inequality reverses:
$g(\E[X]) \ge \E[g(X)]$.
\end{theorem}

\begin{whybox}
The proof is the picture made rigorous with one tool: the \emph{supporting
line}. Because $g$ is convex, at the point $\mu = \E[X]$ there is a line
through $(\mu, g(\mu))$ that stays at or below the whole curve --- a tangent if
$g$ is differentiable, a supporting line in general. Write it as
$L(x) = g(\mu) + m(x - \mu)$, and convexity guarantees $g(x) \ge L(x)$ for all
$x$. Now take expectations of both sides:
\[
  \E[g(X)] \;\ge\; \E[L(X)]
  = g(\mu) + m\,\E[X - \mu]
  = g(\mu) + m\cdot 0
  = g(\E[X]).
\]
The middle term vanished because $\E[X-\mu] = 0$ by definition of $\mu$. The
supporting line, anchored exactly at the mean, is what makes the linear term
disappear and leaves the clean inequality behind.
Figure~\ref{ch19:fig-jensen} shows the curve, a chord, and that supporting line
together.
\end{whybox}

\begin{figure}[H]
\centering
\begin{tikzpicture}
\begin{axis}[
  width=12cm, height=6cm,
  xmin=-0.3, xmax=3.3, ymin=-0.5, ymax=9.5,
  xlabel={$x$}, ylabel={$g(x)$},
  domain=0:3, samples=80,
  legend pos=north west,
  xtick={0.7,1.5,2.3}, xticklabels={$x_1$,$\mu$,$x_2$},
  ytick=\empty,
]
\addplot[densityred, thick] {x^2};
\addlegendentry{convex $g(x)=x^2$}
\addplot[sourceblue, thick] coordinates {(0.7,0.49) (2.3,5.29)};
\addlegendentry{chord (average of $g$)}
\addplot[targetgreen, thick, dashed] {2.25+3*(x-1.5)};
\addlegendentry{supporting line at $\mu$}
\draw[gray, dashed] (axis cs:1.5,0) -- (axis cs:1.5,2.89);
\fill[densityred] (axis cs:1.5,2.25) circle (2pt);
\fill[sourceblue] (axis cs:1.5,2.89) circle (2pt);
\node[below right=2pt, font=\scriptsize, densityred!90!black] at (axis cs:1.5,2.25) {$g(\mu) = g(\E[X])$};
\node[above left=2pt, font=\scriptsize, sourceblue!90!black] at (axis cs:1.5,2.89) {chord height $\E[g(X)]$};
\end{axis}
\end{tikzpicture}
\caption{Jensen's inequality. The convex curve lies below every chord, so the
chord height at $\mu$ --- which equals $\E[g(X)]$ for a two-point variable ---
sits above the curve height $g(\mu) = g(\E[X])$. The dashed supporting line
through $(\mu, g(\mu))$ is the proof: the curve never dips below it, and its
expectation equals $g(\mu)$.}
\label{ch19:fig-jensen}
\end{figure}

\begin{example}[Why variance is non-negative, instantly]
\label{ch19:ex-jensen-var}
Apply Jensen with the convex $g(x) = x^2$:
$\bigl(\E[X]\bigr)^2 \le \E[X^2]$. Rearranged, this is
$\E[X^2] - (\E[X])^2 \ge 0$ --- but the left side is exactly the variance.
Jensen hands you $\Var(X) \ge 0$ in a single line, with equality only when $X$
is a constant (no spread). A fact you proved laboriously before is a one-step
corollary of convexity.
\end{example}

\begin{walkbox}
You drive to a mountain cabin at 30~mph and return along the exact same route
at 60~mph. Is your average speed for the round trip 45~mph, more than 45~mph,
or less than 45~mph?

Before doing any arithmetic, invoke Jensen's inequality! Speed is distance divided
by time, so time per mile is $g(v) = 1/v$. Because $g(v) = 1/v$ is strictly
convex for $v > 0$, the average time per mile is strictly greater than the time
at the average speed:
\[
  \E\!\left[\frac{1}{V}\right] \;>\; \frac{1}{\E[V]}.
\]
Inverting this inequality reverses it: the true average speed is the
\emph{harmonic mean},
$\frac{2}{1/30 + 1/60} = \frac{2}{3/60} = 40\text{ mph}$,
strictly \emph{less} than the naive average of $45\text{ mph}$. You spent
twice as long crawling at 30~mph as you did cruising at 60~mph. Jensen's
inequality predicted that drag before you wrote down a single number.
\end{walkbox}

\begin{paradoxbox}
\textbf{The flaw of averages.} A consultant plans a warehouse. Demand is
random with an average of $100$ units, and profit is some function of demand.
Surely, the consultant reasons, ``planning for average demand gives average
profit.'' This is the \emph{flaw of averages}, and Jensen is the verdict on
it. Whenever profit is a curved function of demand, $\E[g(\text{demand})] \ne
g(\E[\text{demand}])$, and the gap has a guaranteed sign.

A sharp instance: a gambler holds a stock whose value next year is $\$1$ or
$\$4$ with equal probability, average $\$2.50$. If returns \emph{multiply}
over years, the relevant growth is the logarithm, and $\log$ is
\emph{concave}, so Jensen flips: the expected log-growth $\E[\log V]$ is
\emph{strictly less} than $\log \E[V]$. Concretely $\E[\log V] = \tfrac12(\log
1 + \log 4) = \log 2$, while $\log \E[V] = \log 2.5$. Over many years your
wealth compounds at the rate set by $\E[\log V] = \log 2$ (geometric mean
$\$2$), not the arithmetic mean $\$2.50$. Jensen's inequality provides the
mathematical reason the \emph{geometric} mean never exceeds the arithmetic mean.
\end{paradoxbox}

\begin{connectionbox}
Jensen's inequality underpins several major results throughout this book.
It is why \textbf{Kullback--Leibler divergence} $\KL(p\,\|\,q) \ge 0$ (apply
Jensen to the concave $\log$ inside the divergence), a fundamental result in
information theory. It is why the \textbf{EM algorithm} in machine learning
provably increases the likelihood at each step (the E-step builds a Jensen
lower bound on the log-likelihood). And it underlies the variational
inequalities that power modern Bayesian computation. Whenever you see a proof
that ``moves a function inside or outside an expectation'' and gains a definite
direction, look for Jensen hiding underneath.
\end{connectionbox}

\section{A Field Guide to the Inequalities}
\label{ch19:fieldguide}

It helps to see the four tail bounds lined up, because they form a ladder of
increasing assumptions buying increasing strength.

\begin{center}
\begin{tabular}{@{}lll@{}}
\toprule
\textbf{Inequality} & \textbf{Needs} & \textbf{Gives} \\
\midrule
Markov     & $X\ge 0$, mean & $\P(X\ge a) \le \E[X]/a$ \\
Chebyshev  & mean, variance & $\P(\abs{X-\mu}\ge k)\le \sigma^2/k^2$ \\
Chernoff   & MGF, independence & $\le \min_t e^{-ta}M_X(t)$ (exponential) \\
Hoeffding  & bounded, independent & $\le 2e^{-2n\eps^2/(b-a)^2}$ (exponential) \\
\bottomrule
\end{tabular}
\end{center}

Read the table top to bottom as a bargain. Markov asks for almost nothing and
gives a loose $1/a$ bound. Chebyshev adds the variance and earns a $1/k^2$
bound and the Weak Law. Chernoff and Hoeffding add independence and
boundedness and are rewarded with \emph{exponential} decay --- the qualitative
leap that makes large-scale statistics and learning theory possible. The
pattern to internalize is not four separate formulas but one method:
\emph{apply Markov to a cleverly chosen non-negative function}, where each row
just picks a smarter function ($x$, then $x^2$, then $e^{tx}$).

\begin{recoverybox}
If you remember one thing from this chapter, remember the engine: every tail
bound here is Markov's inequality applied to a non-negative function of $X$.
The function determines the strength. Squaring buys variance-based decay;
exponentiating buys exponential decay. Concentration is what happens when you
apply the exponential version to a sum of many independent pieces and watch the
tail collapse like $e^{-cn}$. That collapse is why averages are trustworthy,
why polls work, and why models trained on samples generalize.
\end{recoverybox}

\section{The Bridge to Chapter 20: How Inequalities Power the Laws of Large Numbers}
\label{ch19:bridge}

We began this chapter asking how likely a random quantity is to land far from
its expectation, building from Markov's bound to exponential concentration.
In Chapter 20, we turn this machinery toward infinite horizons to prove the
\textbf{Laws of Large Numbers}. If you average more and more independent trials
forever, does the sample mean $\bar X_n$ converge to the true mean $\mu$? And if
so, in what sense?

The answer splits into two distinct theorems:

\begin{enumerate}
  \item \textbf{The Weak Law of Large Numbers} states that for any fixed tolerance
        $\eps > 0$, the probability of a deviation vanishes in the limit:
        \[
          \lim_{n\to\infty} \P\bigl(\abs{\bar X_n - \mu} \ge \eps\bigr) \;=\; 0.
        \]
        This is \textbf{convergence in probability} (the sample mean lands inside
        the corridor $(\mu-\eps, \mu+\eps)$ with probability approaching $1$).
        You proved this in a single line back in Section~\ref{ch19:chebyshev}:
        Chebyshev's inequality guarantees
        $\P(\abs{\bar X_n - \mu} \ge \eps) \le \frac{\sigma^2}{n\eps^2}$, which
        melts to zero as $n\to\infty$. The $1/n$ rate of Chebyshev is slow, but
        for the Weak Law, all that matters is that it reaches zero.

  \item \textbf{The Strong Law of Large Numbers} makes a vastly more demanding
        claim: with probability $1$, the entire trajectory of running averages
        $\bar X_1, \bar X_2, \bar X_3, \dots$ eventually enters the corridor
        $(\mu - \eps, \mu + \eps)$ and \emph{never leaves it again}. This is
        \textbf{almost sure convergence}.
\end{enumerate}

Why does the Strong Law require more than Chebyshev? The standard bridge to
almost sure convergence is the \textbf{Borel--Cantelli lemma}, which shows
that if the sum of all tail probabilities across time is finite,
\[
  \sum_{n=1}^\infty \P\bigl(\abs{\bar X_n - \mu} \ge \eps\bigr) \;<\; \infty,
\]
then large deviations can happen only a \emph{finite} number of times in total
--- after which the average stays inside the corridor forever.

Here is the key difference. If you feed Chebyshev's bound
$\frac{\sigma^2}{n\eps^2}$ into this sum, the series diverges:
$\sum_{n=1}^\infty \frac{1}{n} = \infty$ (the harmonic series). Chebyshev makes
each individual snapshot probability small, but $1/n$ decays too slowly to rule
out occasional large deviations occurring infinitely often.

Exponential concentration closes that gap. For bounded
variables, Hoeffding's inequality provides
$\P(\abs{\bar X_n - \mu} \ge \eps) \le 2e^{-c n\eps^2}$. Because geometric
series converge, the sum over all time is finite:
\[
  \sum_{n=1}^\infty 2e^{-c n \eps^2} \;<\; \infty.
\]
In one stroke, exponential concentration delivers the summability that the
Borel--Cantelli lemma demands. The collapse of the tails we studied in this
chapter is precisely what allows upgrading from the Weak Law to the Strong Law.
In Chapter 20, we will formalize these modes of convergence and examine how the
sample average settles down.

\section{Exercises}
\label{ch19:exercises}

\begin{exercise}[Markov warm-up]
\label{ch19:ex-markov}
The average score on an exam is $65$ out of $100$. Using only this, give an
upper bound on the fraction of students who scored at least $90$. Then explain
in one sentence why the true fraction is almost certainly far smaller than
your bound, referring to what Markov ignores.
\end{exercise}

\begin{exercise}[Chebyshev in standard units]
\label{ch19:ex-cheb}
A random variable has mean $50$ and standard deviation $4$. Use Chebyshev to
bound $\P(X \le 38 \text{ or } X \ge 62)$. Then find the smallest interval
$[\mu - c\sigma, \mu + c\sigma]$ that Chebyshev guarantees contains at least
$90\%$ of the mass.
\end{exercise}

\begin{exercise}[Sharpness of Markov]
\label{ch19:ex-sharp}
Construct an explicit non-negative random variable for which Markov's
inequality holds with \emph{equality} at $a$, i.e.\ $\P(X\ge a) = \E[X]/a$.
(Hint: a two-point distribution, with mass at $0$ and at $a$.) This shows
Markov cannot be improved in general.
\end{exercise}

\begin{exercise}[Chebyshev to the Weak Law]
\label{ch19:ex-wlln}
Let $X_1,\dots,X_n$ be i.i.d.\ with mean $\mu$ and variance $\sigma^2$. Using
Chebyshev, find the smallest $n$ guaranteeing
$\P(\abs{\bar X_n - \mu} \ge 0.1\sigma) \le 0.01$. Then state precisely what
happens to this probability as $n\to\infty$ and name the theorem.
\end{exercise}

\begin{exercise}[Chernoff for Poisson]
\label{ch19:ex-chernoff-pois}
Let $X \sim \Pois(\lambda)$, whose MGF is $M_X(t) = e^{\lambda(e^t - 1)}$.
Write down the Chernoff bound $\min_{t>0} e^{-ta}M_X(t)$ for $\P(X \ge a)$,
carry out the minimization to find the optimal $t^\star = \ln(a/\lambda)$
(valid for $a>\lambda$), and show the bound simplifies to
$e^{-\lambda}(e\lambda/a)^a$.
\end{exercise}

\begin{exercise}[Hoeffding sample size]
\label{ch19:ex-hoeff-n}
You want to estimate the mean of a quantity known to lie in $[0,10]$ to within
$\eps = 0.5$ with confidence $99\%$. Use the average form of Hoeffding to find
the required number of independent samples $n$. Compare with the count
Chebyshev would demand if you also knew $\sigma = 2$.
\end{exercise}

\begin{exercise}[Jensen direction]
\label{ch19:ex-jensen-dir}
For each function, state whether $\E[g(X)]$ is at least or at most $g(\E[X])$,
and name the convexity that decides it: (a) $g(x) = e^x$; (b)
$g(x) = \sqrt{x}$ on $x \ge 0$; (c) $g(x) = 1/x$ on $x > 0$; (d)
$g(x) = \abs{x}$.
\end{exercise}

\begin{exercise}[AM--GM from Jensen]
\label{ch19:ex-amgm}
Let $X$ take the values $x_1,\dots,x_n > 0$ each with probability $1/n$. Apply
Jensen to the concave function $g(x) = \log x$ and exponentiate the result to
prove the arithmetic mean--geometric mean inequality
$\frac1n\sum x_i \ge \bigl(\prod x_i\bigr)^{1/n}$. State the condition for
equality.
\end{exercise}

\begin{exercise}[Why the exponential beats the square --- a proof]
\label{ch19:ex-proof}
Let $S_n$ be the number of heads in $n$ fair flips. (a) Use Chebyshev to bound
$\P(S_n \ge 0.6n)$ and show the bound decays like $C/n$. (b) Use the Chernoff
result $\P(S_n \ge (\tfrac12+\delta)n) \le e^{-2\delta^2 n}$ to bound the same
event. (c) For $n = 500$, compute both numerically and explain in two or three
sentences, in terms of the optimization over $t$, why the exponential method
captures a decay the quadratic method structurally cannot.
\end{exercise}

\begin{exercise}[Conceptual: what concentration is not]
\label{ch19:ex-concept}
A student claims: ``Concentration means that if I flip a coin many times, I am
guaranteed to get exactly half heads.'' Explain, using the precise statement of
Hoeffding's inequality, in what sense this is wrong and in what sense it is
nearly right. Distinguish carefully between the \emph{value} concentrating and
the randomness \emph{disappearing}.
\end{exercise}

\begin{exercise}[Summability and the bridge to the Strong Law]
\label{ch19:ex-summability}
Let $X_1, X_2, \dots$ be independent random variables bounded in $[0,1]$ with
mean $\mu$, and let $\bar X_n = \frac1n \sum_{i=1}^n X_i$. Fix $\eps > 0$.
\begin{enumerate}[label=(\alph*)]
  \item Apply Hoeffding's inequality to write an upper bound for the tail
        probability $p_n = \P(\abs{\bar X_n - \mu} \ge \eps)$.
  \item Compute the infinite sum $\sum_{n=1}^\infty p_n$ by bounding it with
        a geometric series. Show that this sum converges to a finite value.
  \item Contrast this with the bound that Chebyshev's inequality gives:
        $\P(\abs{\bar X_n - \mu} \ge \eps) \le \frac{1}{4n\eps^2}$. What is
        $\sum_{n=1}^\infty \frac{1}{4n\eps^2}$? Explain in two sentences why
        Chebyshev proves the Weak Law of Large Numbers but cannot prove the
        Strong Law via the Borel--Cantelli lemma.
\end{enumerate}
\end{exercise}
